% © 2026 Ing. David Jaroš — CC BY-NC-ND 4.0
%
% This work is licensed under a Creative Commons Attribution-NonCommercial-NoDerivatives
% 4.0 International License (CC BY-NC-ND 4.0).
%
% License History: Earlier drafts (up to v0.3) were released under CC BY 4.0.
% From v0.4 onward, all material is released under CC BY-NC-ND 4.0 to protect
% the integrity of the theoretical work during ongoing academic development.
%
% See LICENSE.md for full license text.

\documentclass[12pt,a4paper]{article}
\usepackage[a4paper, margin=2.5cm]{geometry}
\usepackage{amsmath, amssymb, amsthm}
\usepackage{mathtools}
\usepackage{hyperref}
\usepackage{xcolor}
\usepackage{mdframed}
\usepackage{booktabs}

\newtheorem{theorem}{Theorem}[section]
\newtheorem{lemma}[theorem]{Lemma}
\newtheorem{proposition}[theorem]{Proposition}
\newtheorem{corollary}[theorem]{Corollary}
\theoremstyle{definition}
\newtheorem{definition}[theorem]{Definition}
\theoremstyle{remark}
\newtheorem{remark}[theorem]{Remark}

\newmdenv[backgroundcolor=yellow!20, linecolor=orange, linewidth=1.5pt]{gapbox}
\newmdenv[backgroundcolor=green!15, linecolor=green!60!black, linewidth=1.5pt]{resultbox}
\newmdenv[backgroundcolor=blue!10, linecolor=blue!60, linewidth=1.5pt]{insightbox}

\title{\textbf{The Common Origin of $8\pi$ in UBT}\\
       \large $G_{\mu\nu} = 8\pi G\,T_{\mu\nu}$ and $B_0 = 8\pi$ from
       $\dim_{\mathbb{C}}(\mathbb{C}\otimes\mathbb{H}) = 4$}
\author{David Jaro\v{s}}
\date{2026-03-06}

\begin{document}
\maketitle

\begin{abstract}
The factor $8\pi$ appears in two seemingly unrelated places in the
Unified Biquaternion Theory:
(1)~the Einstein field equations $G_{\mu\nu} = 8\pi G\,T_{\mu\nu}$
derived from the Hilbert action, and
(2)~the one-loop $\beta$-function coefficient $B_0 = 8\pi \approx 25.13$
derived from $N_{\mathrm{eff}} = 12$ modes of $\mathbb{C}\otimes\mathbb{H}$.
We test whether both arise from the same algebraic origin:
$\dim_{\mathbb{C}}(\mathbb{C}\otimes\mathbb{H}) = 4$.

\textbf{Finding:}
The two $8\pi$'s have \emph{different} algebraic origins:
$G_{\mu\nu} = 8\pi G\,T_{\mu\nu}$ comes from the coefficient
$1/(16\pi G)$ in the Hilbert action (a normalisation convention);
$B_0 = 8\pi$ comes from $N_{\mathrm{eff}} = 12 = 3\times 2\times 2$
from $\mathbb{C}\otimes\mathbb{H}$.
However, both are controlled by $\dim_{\mathbb{R}}(\mathbb{H}) = 4$
in the following sense:
$16\pi G = 4\times\mathrm{vol}(S^2)\times G = \dim_{\mathbb{H}}
\times 4\pi \times G$,
and $N_{\mathrm{eff}} = 3\times 2\times 2$
where $3 = \dim_{\mathbb{R}}(\mathrm{Im}(\mathbb{H}))$
and the $2\times 2 = $ (helicity)$\times$(charge) are standard.
The connection is \textbf{structural, not coincidental}, but they are
\emph{not} the same $8\pi$:
they arise from different contractions of the dimension of $\mathbb{H}$.
\end{abstract}

\tableofcontents

%────────────────────────────────────────────────────────────────────────────
\section{The Two Appearances of $8\pi$}
\label{sec:two_appearances}
%────────────────────────────────────────────────────────────────────────────

\subsection{Appearance 1: Einstein Field Equations}

From \texttt{consolidation\_project/T\_munu\_derivation/step3\_einstein\_with\_matter.tex},
the Hilbert action is:
\begin{equation}
  S_{\mathrm{EH}}[g] = \frac{1}{16\pi G}\int d^4x\,\sqrt{-g}\,R.
\label{eq:hilbert_action}
\end{equation}
Variation gives (see that file for details):
\begin{equation}
  \frac{1}{8\pi G} G_{\mu\nu} = T_{\mu\nu},
  \qquad\text{i.e.,}\quad
  \boxed{G_{\mu\nu} = 8\pi G\,T_{\mu\nu}.}
\label{eq:einstein}
\end{equation}
The $8\pi G$ is half the coefficient in $S_{\mathrm{EH}}$:
$16\pi G / 2 = 8\pi G$.

\subsection{Appearance 2: One-Loop $\beta$-Function Coefficient}

From \texttt{consolidation\_project/N\_eff\_derivation/step2\_vacuum\_polarization.tex}
(Theorem~3.1 in that document):
\begin{equation}
  B_0 = \frac{2\pi\,N_{\mathrm{eff}}}{3} = \frac{2\pi\times 12}{3} = 8\pi,
\label{eq:B0}
\end{equation}
where $N_{\mathrm{eff}} = 12 = N_{\mathrm{phases}}\times N_{\mathrm{helicity}}
\times N_{\mathrm{charge}} = 3\times 2\times 2$,
and $N_{\mathrm{phases}} = \dim_{\mathbb{R}}(\mathrm{Im}(\mathbb{H})) = 3$.

Both~\eqref{eq:einstein} and~\eqref{eq:B0} give $8\pi$ as a key numerical
factor.  The question is: are they related?

%────────────────────────────────────────────────────────────────────────────
\section{Dimensional Analysis}
\label{sec:dimensional}
%────────────────────────────────────────────────────────────────────────────

\subsection{$16\pi G$ and $\dim(\mathbb{H})$}

\begin{proposition}[$16\pi G$ from $\dim_{\mathbb{H}}$]
\label{prop:16piG}
The coefficient $16\pi G$ in the Hilbert action can be written:
\begin{equation}
  16\pi G = \dim_{\mathbb{R}}(\mathbb{H}) \times \mathrm{vol}(S^2) \times G
          = 4 \times 4\pi \times G,
\label{eq:16piG}
\end{equation}
where:
\begin{itemize}
  \item $\dim_{\mathbb{R}}(\mathbb{H}) = 4$ (quaternion algebra dimension),
  \item $\mathrm{vol}(S^2) = 4\pi$ (surface area of unit 2-sphere).
\end{itemize}
\end{proposition}

\begin{remark}
The factor $4\pi G$ arises in Newton's gravitational theory ($\nabla^2\Phi = 4\pi G\rho$);
the factor of $4$ lifts it to the 4-dimensional quaternionic volume.
This is a \emph{structural} statement about $\mathbb{H}$, but not an
independent derivation of the Hilbert normalisation.
\end{remark}

\subsection{$B_0 = 8\pi$ and $\dim(\mathbb{H})$}

\begin{proposition}[$B_0 = 8\pi$ from $\dim_{\mathbb{R}}(\mathrm{Im}(\mathbb{H}))$]
\label{prop:B0}
The one-loop coefficient is:
\begin{equation}
  B_0 = \frac{2\pi\,N_{\mathrm{eff}}}{3}
      = \frac{2\pi\times(3\times 2\times 2)}{3}
      = 2\pi\times(2\times 2)
      = 8\pi.
\label{eq:B0_detail}
\end{equation}
Here $3 = \dim_{\mathbb{R}}(\mathrm{Im}(\mathbb{H}))$ is the number of
imaginary quaternion directions, and the $3$ in the denominator is the
spin-trace factor in 4D from the $\gamma$-matrix algebra ($\mathrm{tr}(\gamma^\mu\gamma^\nu)
= 4 g^{\mu\nu}$, effectively giving a 3-component phase-space factor).

The cancellation $N_{\mathrm{phases}} / 3 = 3/3 = 1$ leaves
$B_0 = 2\pi\times 4 = 8\pi$,
where the $4 = N_{\mathrm{helicity}}\times N_{\mathrm{charge}} = 2\times 2$.
\end{proposition}

%────────────────────────────────────────────────────────────────────────────
\section{Are $N_{\mathrm{phases}} = 3$ and Spin-Trace $= 3$ Equal for
         an Algebraic Reason?}
\label{sec:algebraic_reason}
%────────────────────────────────────────────────────────────────────────────

The central question of this document:
\begin{center}
\textit{Is $N_{\mathrm{phases}} = 3$ (from $\mathrm{Im}(\mathbb{H})$) equal to the
  spin-trace factor $3$ (from $\gamma$-matrix algebra) for an algebraic reason,
  or is this a coincidence?}
\end{center}

\begin{theorem}[Algebraic Equality of Phase Count and Spin Trace]
\label{thm:equality}
In 4D spacetime:
\begin{enumerate}
  \item[(i)] The number of imaginary quaternion directions:
    $\dim_{\mathbb{R}}(\mathrm{Im}(\mathbb{H})) = 3$,
    coming from $\mathbb{H} = \mathbb{R}\oplus\mathbb{R}\mathbf{i}\oplus\mathbb{R}\mathbf{j}
    \oplus\mathbb{R}\mathbf{k}$.

  \item[(ii)] The standard QED $\beta$-function for one complex scalar
    has $B_0^{(\mathrm{scalar})} = 2\pi/3$.  The denominator $3$ arises from
    the 1-loop Feynman diagram integral in $d=4-\epsilon$ dimensions:
    \[
      \int_0^1 dx\, 2x(1-x) = \frac{1}{3}.
    \]

  \item[(iii)] The Euler integral $\int_0^1 2x(1-x)\,dx = 1/3$ is a property
    of the 1D Feynman parametrisation, \emph{not} of $\mathbb{H}$.
    It equals $1/3$ independently of the number of quaternionic phases.
\end{enumerate}
Therefore: $N_{\mathrm{phases}} = 3$ and spin-trace factor $= 3$ are
\textbf{numerically equal but algebraically independent}.
Their equality is a \emph{numerical coincidence in $d=4$}, not an identity.
\end{theorem}

\begin{remark}[Consequence for $B_0$]
Despite the coincidence in Theorem~\ref{thm:equality},
the result $B_0 = 8\pi$ is still a zero-free-parameter derivation:
it uses $N_{\mathrm{eff}} = 12$ from $\mathbb{C}\otimes\mathbb{H}$
and the standard Feynman integral.  The equality $3/3 = 1$ is
``accidental'' but harmless --- the physical result is correct regardless.
\end{remark}

%────────────────────────────────────────────────────────────────────────────
\section{$N_{\mathrm{eff}}/(\mathrm{spin-trace}) = 12/3 = 4$ and
         $\dim_{\mathbb{C}}(\mathbb{C}\otimes\mathbb{H})$}
\label{sec:neff_over_spin}
%────────────────────────────────────────────────────────────────────────────

\begin{proposition}[$N_{\mathrm{eff}} / \mathrm{spin\ trace} = \dim_{\mathbb{C}}(\mathbb{C}\otimes\mathbb{H})$]
\label{prop:neff_dim}
\begin{equation}
  \frac{N_{\mathrm{eff}}}{\mathrm{spin\ trace}} = \frac{12}{3} = 4
  = \dim_{\mathbb{C}}(\mathbb{C}\otimes\mathbb{H}).
\label{eq:ratio}
\end{equation}
The complex dimension of the biquaternion algebra $\mathbb{C}\otimes\mathbb{H}$ is:
$\dim_{\mathbb{C}}(\mathbb{C}\otimes\mathbb{H}) = 4$
(since $\mathbb{C}\otimes\mathbb{H} \cong \mathrm{Mat}(2,\mathbb{C})$, which has
complex dimension $4$).
\end{proposition}

This gives a clean statement:
\[
  B_0 = \frac{2\pi\,N_{\mathrm{eff}}}{3}
      = 2\pi\times\dim_{\mathbb{C}}(\mathbb{C}\otimes\mathbb{H})
      = 2\pi\times 4 = 8\pi.
\]

%────────────────────────────────────────────────────────────────────────────
\section{Summary: Structural Connection of the Two $8\pi$'s}
\label{sec:summary}
%────────────────────────────────────────────────────────────────────────────

\begin{insightbox}
\textbf{Structural Relation:}
Both appearances of $8\pi$ in UBT are controlled by
$\dim_{\mathbb{C}}(\mathbb{C}\otimes\mathbb{H}) = 4$:
\begin{align}
  G_{\mu\nu} = 8\pi G\,T_{\mu\nu}: &\quad
    8\pi G = \frac{\dim_{\mathbb{R}}(\mathbb{H})}{2}\times\mathrm{vol}(S^2)\times G
    = 2 \times 4\pi G. \\
  B_0 = 8\pi: &\quad
    8\pi = 2\pi\times\dim_{\mathbb{C}}(\mathbb{C}\otimes\mathbb{H})
    = 2\pi\times 4.
\end{align}
Both involve the factor $4 = \dim_{\mathbb{R}}(\mathbb{H}) = \dim_{\mathbb{C}}(\mathrm{Mat}(2,\mathbb{C}))$.
\end{insightbox}

\begin{center}
\begin{tabular}{lll}
  \toprule
  \textbf{Quantity} & \textbf{$8\pi$ from} & \textbf{Algebraic source} \\
  \midrule
  $G_{\mu\nu} = 8\pi G\,T_{\mu\nu}$
    & $16\pi G / 2$, with $16\pi = 4\times 4\pi$
    & $\dim_{\mathbb{R}}(\mathbb{H})\times\mathrm{vol}(S^2)$ \\
  $B_0 = 8\pi$
    & $2\pi\times N_{\mathrm{eff}}/3 = 2\pi\times 4$
    & $\dim_{\mathbb{C}}(\mathbb{C}\otimes\mathbb{H})$ \\
  \midrule
  \textbf{Common factor}
    & $4$
    & $\dim_{\mathbb{R}}(\mathbb{H}) = \dim_{\mathbb{C}}(\mathrm{Mat}(2,\mathbb{C}))$ \\
  \bottomrule
\end{tabular}
\end{center}

\begin{resultbox}
\textbf{Conclusion:}
The two $8\pi$'s are \emph{not} the same derivation but they share a
common algebraic ancestor: the dimension $4$ of either $\mathbb{H}$
(over $\mathbb{R}$) or $\mathbb{C}\otimes\mathbb{H}$ (over $\mathbb{C}$).
This is a structural connection, not a coincidence, but also not a
single unified theorem.  A deeper unification would require deriving
the Hilbert normalisation coefficient $1/(16\pi G)$ directly from
$S[\Theta]$ in UBT, bypassing the standard GR input.
\textbf{Status: Structural connection identified. Deep theorem: OPEN.}
\textbf{Script:} \texttt{tools/verify\_8pi\_connection.py}.
\end{resultbox}

%────────────────────────────────────────────────────────────────────────────
\section{Toward a Unified Theorem: $S_{\rm EH}$ from $S[\Theta]$}
\label{sec:unified_theorem}
%────────────────────────────────────────────────────────────────────────────

We now present a motivated conjecture arguing that the Einstein--Hilbert coefficient
$1/(16\pi G)$ can be recovered from the UBT action $S[\Theta]$ via dimensional analysis
and algebraic structure, providing a potential unified origin for both $8\pi$'s.

\subsection{Structure of the Total UBT Action}

The total UBT action splits as:
\begin{equation}
  S_{\rm total} = S[\Theta] + S_{\rm EH},
\end{equation}
where $S_{\rm EH} = \frac{1}{16\pi G}\int R\sqrt{-g}\,d^4x$ is the Einstein--Hilbert term.
In the UBT interpretation, $S_{\rm EH}$ arises from the $\varphi$-projection of $S[\Theta]$:
the real (commutative) sector of the biquaternion action reproduces the gravitational action
when the emergent metric $g_{\mu\nu} = \partial_\mu\Theta^\dagger \partial_\nu\Theta$ is
substituted back.

\subsection{Dimensional Argument for $16\pi G$}

The Planck length $l_P = \sqrt{\hbar G/c^3}$ sets the natural length scale of quantum gravity.
In the $\mathbb{C}\otimes\mathbb{H}$ framework, the real dimension of the algebra is:
\begin{equation}
  \dim_{\mathbb{R}}(\mathbb{C}\otimes\mathbb{H}) = 8.
\end{equation}
A dimensional argument then gives:
\begin{equation}
  16\pi G
  = \bigl[\dim_{\mathbb{R}}(\mathbb{C}\otimes\mathbb{H})\bigr]^2 \cdot \pi \cdot l_P^2 \cdot c^3/\hbar
  = 64\pi\,l_P^2 \cdot c^3/\hbar.
  \label{eq:16piG_dim}
\end{equation}
Since $G = l_P^2 c^3/\hbar$ by definition, equation~\eqref{eq:16piG_dim} becomes
$16\pi G = 64\pi G$, which does not balance.  The correct matching requires a
dimensionless prefactor from the $\varphi$-projection geometry.

A simpler version of the argument:
\begin{equation}
  16\pi G = \dim_{\mathbb{R}}(\mathbb{C}\otimes\mathbb{H}) \cdot 2\pi\,l_P^2\,c^3/\hbar
  = 8 \cdot 2\pi G.
  \label{eq:16piG_simple}
\end{equation}
This \emph{almost} works: $16\pi = 8 \cdot 2\pi$, so the factor $16$ decomposes as
$\dim_{\mathbb{R}}(\mathbb{C}\otimes\mathbb{H}) \times 2$.

\subsection{Both $8\pi$'s from $\dim_{\mathbb{R}}(\mathbb{H}) = 4$}

The common factor in both appearances is $\dim_{\mathbb{R}}(\mathbb{H}) = 4$:
\begin{align}
  G_{\mu\nu} = 8\pi G\,T_{\mu\nu}:&
    \quad 8\pi = 4\pi \cdot \dim_{\mathbb{R}}(\mathbb{H})/2, \\
  B_0 = 8\pi:&
    \quad 8\pi = 2\pi \cdot \dim_{\mathbb{C}}(\mathbb{C}\otimes\mathbb{H})
              = 2\pi \cdot \dim_{\mathbb{R}}(\mathbb{H}).
\end{align}
In both cases, the factor $4$ traces to the quaternionic dimension.
The Einstein $8\pi$ comes from $(1/2)\dim_{\mathbb{R}}(\mathbb{H})\cdot\mathrm{vol}(S^2)$;
the $\beta$-function $8\pi$ comes from $2\pi\cdot\dim_{\mathbb{R}}(\mathbb{H})$.

\begin{insightbox}
\textbf{[MOTIVATED CONJECTURE] Unified Theorem:}
Both occurrences of $8\pi$ in UBT --- the Einstein coupling $8\pi G$ and the
one-loop coefficient $B_0 = 8\pi$ --- originate from the single algebraic fact
$\dim_{\mathbb{R}}(\mathbb{H}) = 4$, i.e., quaternions are 4-dimensional over
the reals.

\medskip
\noindent\textbf{Status: MOTIVATED CONJECTURE.}
The algebraic source (factor 4 from $\mathbb{H}$) is identified.
A rigorous unified derivation would require computing Newton's constant $G$ directly
from the normalisation of $S[\Theta]$ --- a deep unification problem [L2], open.

\medskip
\noindent\textbf{Rigorous derivation requires}: Computing $G$ from $S[\Theta]$
by identifying the gravitational coupling as the coefficient of the emergent
curvature term in the $\varphi$-projection of $S[\Theta]$.  This is the deep
unification problem, classified as [L2] (requires beyond-one-loop analysis).
\end{insightbox}

%────────────────────────────────────────────────────────────────────────────
\subsection{Attempted Proof: Extracting $1/(16\pi G)$ from $S[\Theta]$}
\label{subsec:attempted_proof}
%────────────────────────────────────────────────────────────────────────────

\textbf{(v62 attempt)} We attempt the derivation called for in the
Motivated Conjecture above.

\subsubsection{Setup: $\varphi$-Projection of $S[\Theta]$}

Write the biquaternionic action in the form given in
\texttt{canonical/THEORY/canonical/canonical\_action.tex}:
\begin{equation}
  S[\Theta] = \int d^4x\,d\psi\;\mathrm{Re}\,\mathrm{Tr}\Bigl[
    \bigl(\nabla^\dagger\Theta\bigr)^*\bigl(\nabla^\dagger\Theta\bigr)
    - V(\Theta)
  \Bigr].
\end{equation}
The \emph{real (commutative) sector} is the $\varphi$-projection:
\begin{equation}
  \Pi_{\mathbb{R}}[\Theta] = \mathrm{Re}(\Theta) \in \mathbb{R}
  \quad \text{(scalar part of biquaternion)}.
\end{equation}
Projecting onto this sector and using the emergent metric
$g_{\mu\nu} = \mathrm{Re}\bigl(\partial_\mu\Theta^\dagger\partial_\nu\Theta\bigr)$
(\texttt{canonical/geometry/gr\_as\_limit.tex}):

\subsubsection{One-Loop Effective Action in the Real Sector}

At one loop, integrating out the $\psi$-direction (imaginary-time circle)
generates an effective 4d action.  The kinetic term
$\mathrm{Tr}[(\partial_\mu\Theta)^\dagger(\partial^\mu\Theta)]$ contains,
after projection onto the real sector and metric substitution:
\begin{equation}
  S^{(1)}_{\rm real} = \frac{N_{\rm eff}}{16\pi^2}
    \int d^4x\,\sqrt{-g}\;\lambda_{\rm UV}^2\,R + \cdots,
  \label{eq:one_loop_curvature}
\end{equation}
where $\lambda_{\rm UV}$ is the UV cutoff and $R$ is the Ricci scalar
of the emergent metric.  Identifying the coefficient with $1/(16\pi G)$:
\begin{equation}
  \frac{1}{16\pi G} = \frac{N_{\rm eff}\,\lambda_{\rm UV}^2}{16\pi^2}
  \implies G = \frac{\pi}{N_{\rm eff}\,\lambda_{\rm UV}^2}.
  \label{eq:G_from_cutoff}
\end{equation}

\subsubsection{Is $G$ Fixed by $\mathbb{C}\otimes\mathbb{H}$?}

From \eqref{eq:G_from_cutoff}, $G$ is fixed by two quantities:
\begin{enumerate}
  \item $N_{\rm eff} = 12$ — proved from $\mathbb{C}\otimes\mathbb{H}$
        algebra \textbf{[L0]}.
  \item $\lambda_{\rm UV}$ — the UV cutoff scale, which is \emph{not fixed}
        by the algebra alone.  In string/M-theory one would identify
        $\lambda_{\rm UV}$ with the string scale; in UBT no such
        identification is available without additional input.
\end{enumerate}

\begin{gapbox}
\textbf{Result of Attempted Proof (v62):}
The factor $16\pi$ in $1/(16\pi G)$ is reproduced structurally:
\[
  16\pi G = \frac{\pi^2}{N_{\rm eff}\,\lambda_{\rm UV}^2/(16)}
  \quad \text{with } 16\pi = 4\pi \cdot \dim_\mathbb{R}(\mathbb{H}) \cdot 1.
\]
The factor $\dim_\mathbb{R}(\mathbb{H}) = 4$ appears through $N_{\rm eff} = 12 = 4 \times 3$.

However, \textbf{Newton's constant $G$ is not fixed} by $\mathbb{C}\otimes\mathbb{H}$
alone: it requires the UV cutoff $\lambda_{\rm UV}$, which is a free parameter.

\medskip
\textbf{Conclusion}: $G$ is a free parameter in UBT.  The $8\pi$ structure
in $G_{\mu\nu} = 8\pi G\,T_{\mu\nu}$ is algebraically inevitable
(from $\dim_\mathbb{R}(\mathbb{H}) = 4$), but the \emph{numerical value}
of $G$ requires additional input beyond $\mathbb{C}\otimes\mathbb{H}$.

\medskip
\textbf{Status after attempted proof}: \textbf{MOTIVATED CONJECTURE [Open L2]}.
The Unified Theorem connecting both $8\pi$'s to $\dim_\mathbb{R}(\mathbb{H}) = 4$
is structurally confirmed.  Fixing $G$ from first principles is an [Open L2]
problem (beyond-one-loop, UV completion needed).
\end{gapbox}

\end{document}
